By definition, the term apocalypse is defined as a cataclysmic event, which is
from the etymonline dictionary. The events which Primo faced can be described as
an apocalypse, or the end of a world and the creation of a new world, which is
an appropriate term to use when describing the \textbf{Story of Ten Days}. Primo
faces unimaginable events.

In the middle of the story \textbf{The Story of Ten Days}, Primo describes the
Lager as \textit{"The Lager, hardly dead, had already begun to decompose \ldots"
(Pg. $158$)} Eventually, in the story, we see that the decomposition of the camp
is essential in creating a new world for Primo and the people with him.
\textit{"The work of the bombs had been completed by the work of man \ldots (Pg.
$158$)"} This statement, appears to be criticising or telling that man is
responsible for the destruction of their own kind. Then Primo goes on to
describe the people who are suffering in the camp. \textit{"\ldots ragged,
decrepit, skeleton- like patients at all able to move dragged themselves
everywhere on the frozen soil, like an invasion of worms \ldots (Pg. $158$)"}
Here, Primo is comparing the people who are suffering in the camps to worms.
This may indicate that he's telling us that the Lager has transformed the people
into the lowest of life form, worms. \textit{"no longer in control of their own
bowels, they had fouled everywhere, polluting the precious snow, the only source
of water remaining in the whole camp \ldots" (Pg. $158$)} Primo describes the suffering
people with disgust, but when he mentions the snow, he elevates its status by
calling precious, which indicates that the snow, to him, is more valuable than
those people, which you cannot blame him. Without the snow, he wouldn't have any
water to drink and those people keep depleting the amount of water they have 

% \textit{"The work of the bombs had been completed by the work of man"}. This
% statement, appears to be criticising or telling that man is responsible for the
% destruction of their own kind. \textit{"ragged, decrepit, skeleton-like patients
% at all able to move dragged themselves everywhere on the frozen soil, like an
% invasion of worms"}. Primo describes the people in a de-humanizing way, and
% compares them to an invasion of worms, which may be how he sees them. Just a
% bunch of worms who are going to eventually parish. \textit{"no longer in control
% of their own bowels, they had fouled everywhere, polluting the precious snow,
% the only source of water remaining in the whole camp."} Primo describes those
% people with disgust, but when he mentions the snow, he elevates its status by
% calling \textbf{precious}, which indicates that the snow, to him, is more
% valuable than those people, which you cannot blame him. Without the snow, he
% wouldn't have any water to drink and those people keep depleting the amount of
% water they have. This shows the affect the camp has on Primo. At this point,
% Primo is starting to lose the humanity that he has.

% You can tell the exact moment where Primo and everybody else have given up on
% humanity. The law of the Lager is: \textit{"eat your own bread, and if you can,
% that of your neighbour"}. But there is still a sense of unity, which is
% indicated by: \textit{"I had my bed here; and by now a tie united us, the eleven
% patients of the Infektionsabteilung"}. Even when they are losing all hope, all
% the humanity in them, there is still a sense of community, which is something
% vital when in times of desperation.

% Very occasionally we heard the thundering of artillery, both near and far, and
% at intervals, the crackling of automatic rifles. In the darkness, lighted only
% by the glow of the embers, Arthur and I sat smoking cigarettes made of herbs
% found in the kitchen and spoke of many things, both past and future. In the
% middle of this endless plain, frozen and full of war, in the small dark room
% swarming with germs, we felt at peace with ourselves and with the world. We
% were broken by tiredness, but we seemed to have finally accomplished something
% useful - perhaps like God after the first day of creation.

% I was thinking life outside was beautiful and would be beautiful again, and
% that it would really be a pity to let ourselves be overcome now. I woke up the
% patients who were dozing and when I was sure that they were all listening I
% told them, first in French and then in my best German, that they must all
% begin to think of returning home now, and that as far as depended on us,
% certain things were to be done and others to be avoided. Each person should
% carefully look after his own bowl and spoon; no one should offer his own soup
% to others; no one should climb down from his bed except to go to the latrine;
% if anyone was in need of anything, he should only turn to us three. Arthur in
% particular was given the task of supervising the discipline and hygiene, and
% was to remember that it was better to leave bowls and spoons dirty rather than
% wash them with the danger of changing those of a diphtheria patient with those
% of someone suffering from typhus.
